\subsection{A reduced counterexample with nonfinite normalization}

\begin{lemma}[A differential subring with an analytically ramified completion]
\label{lem:C-RD11-ring}
Let $V$ be a complete discrete valuation ring of mixed characteristic
$(0,p)$ with uniformizer $p$, fraction field $K$ and residue field $k$.
Let $C\subset K$ be a subfield such that $C\cap V\to k$ is onto.
Let $J\in V$ be transcendental over $C$, let $F\subset K$ be an
algebraic extension of $C(J)$, and let $D:F\to F$ extend the derivation
which is zero on $C$ and sends $J$ to $1$. Put
\[
 V_0=F\cap V,\qquad R=\{a\in V_0:D(a)\in V_0\}.
\]
Then $R$ is a one-dimensional noetherian local domain, with fraction
field $F$ and residue field $k$. Its normalization is the DVR $V_0$,
which has uniformizer $p$ and residue field $k$; both the fraction-field
and residue-field extensions induced by $R\subset V_0$ are identities.
In particular $R$ is geometrically unibranch. Its normalization is not
finite, and the homomorphism
\[
 R\longrightarrow V[\epsilon]/(\epsilon^2),
 \qquad a\longmapsto a+\epsilon D(a)
                                                        \tag{C-RD11.1}
\]
identifies its maximal-ideal completion with
$V[\epsilon]/(\epsilon^2)$.
\end{lemma}
\begin{proof}
The ring $V_0$ is a valuation ring with value group $\mathbf Z$ and
uniformizer $p$, hence is a DVR. Every $f\in F$ can be multiplied by
a sufficiently large power of $p$ to make both $f$ and $D(f)$ integral;
therefore $F=R[1/p]=\operatorname{Frac}R$. Elements of $R$ of valuation
zero are units of $R$, since $D(a^{-1})=-a^{-2}D(a)$ is integral.
Thus $R$ is local, its maximal ideal is
$\mathfrak m=\{a\in R:v(a)>0\}$, and its residue field is $k$.

Choose $c\in C\cap V$ with the same residue as $J$ and put $z=J-c$.
Then $z\in\mathfrak m$ and $D(z)=1$. For $a\in\mathfrak m$, choose
$b\in C\cap V$ lifting the residue of $D(a)$. Both $a-bz$ and its
derivative are divisible by $p$, so $(a-bz)/p\in R$. Consequently
\[
 \mathfrak m=(p,z),\qquad \mathfrak m^2=p\mathfrak m,
 \qquad \mathfrak m^n=p^{n-1}\mathfrak m\quad(n\geq1). \tag{C-RD11.2}
\]
For the middle equality, observe that $z^2/p\in\mathfrak m$: its
valuation is positive and its derivative $2z/p$ is integral.

If $0\ne a\in R$ has valuation $r\geq1$ and
$b\in\mathfrak m^{2r}$, then
$v(b)\geq2r$, $v(D(b))\geq2r-1$ and $v(D(a))\geq0$. Hence
$b/a$ and
\[
 D(b/a)=\frac{aD(b)-bD(a)}{a^2}
\]
are integral. Thus $\mathfrak m^{2r}\subset aR$.
Since $\mathfrak m$ is finitely generated and $R/\mathfrak m=k$,
$R/aR$ is an artinian noetherian ring. Every nonzero ideal $I$ contains
such an $a$ and is finitely generated by lifting generators of $I/aR$.
This proves that $R$ is noetherian. Every nonzero prime contains a
power of $\mathfrak m$, so $\dim R=1$.

The residue-surjectivity assumption makes $C\cap V$ dense in $V$:
successively choose representatives of the residues of the remainders
after division by $p$. For $a\in V_0$, choose $N$ sufficiently large
that $p^Na^{p^N-1}D(a)$ is integral. Then $a^{p^N}\in R$, so $a$
is integral over $R$.
As $V_0$ is integrally closed, it is exactly the normalization. The
normalization has one point above the closed point and the same
residue field, which is the geometric-unibranch criterion.

Write $T=V[\epsilon]/(\epsilon^2)$ and
$\mathfrak n=(p,\epsilon)$. Formula (C-RD11.2) shows that the
inverse image of $\mathfrak n^n=p^nV+\epsilon p^{n-1}V$ under
(C-RD11.1) is exactly $\mathfrak m^n$. The image is dense: given
$a,b\in V$, approximate them by $a_0,b_0\in C\cap V$, and choose
$c_N\in C\cap V$ with $v(J-c_N)\geq N$. Then
\[
 r=a_0+b_0(J-c_N)\in R,\qquad D(r)=b_0,
\]
approximates $a+\epsilon b$ to the required order. Thus
$R/\mathfrak m^n\simeq T/\mathfrak n^n$ for every $n$, proving the
completion assertion.

Finally no nonzero $d\in R$ can satisfy $dV_0\subset R$. With $c_N$
as above put $u_N=(J-c_N)/p^N\in V_0$. Then
\[
 D(du_N)=D(d)u_N+d/p^N.
\]
For $N>v(d)$ the first term is integral and the second has negative
valuation, so $du_N\notin R$. A finite birational normalization
would have a nonzero conductor, by taking a common denominator for
finitely many generators. Therefore $V_0/R$ is not finite.
\end{proof}

\begin{theorem}[C-RD11: the reduced assertion fails without finite normalization]
\label{thm:C-RD11}
\status{Proved; independent audit of the construction recorded.}
There is a reduced, one-dimensional, geometrically unibranch,
noetherian local domain $R$ of mixed characteristic $(0,5)$, with
nonfinite normalization, for which hypotheses a) and b) of 100:3
hold, but the functor $F$ is not essentially surjective.
The failure already occurs for a split ordinary height-two BT group
and an elliptic generic fibre which itself has a good model over $R$.
Thus the assertion in C.2 for every reduced base is false.
\end{theorem}
\begin{proof}
Let $k=\overline{\mathbf F_5(t)}$ and $V=W(k)$.
Take the elliptic curve over $k$ with $j=t$. It is ordinary: in the
short model
\[
 Y^2=X^3-\frac{3t}{t-1728}X-\frac{2t}{t-1728},
\]
the coefficient of $X^4$ in the square of the cubic, namely twice
its $X$ coefficient, is nonzero in characteristic $5$.
Let $E^{\mathrm{can}}/V$ be its canonical ordinary lift and write
$J=j(E^{\mathrm{can}})$. The canonical split BT description is
\cite[Theorem~2.1]{Katz1981SerreTate}; in relative dimension one the
formal lift algebraizes using the relatively ample divisor given by
three times the zero section. We have
\[
 E^{\mathrm{can}}[5^\infty]\simeq
 \mu_{5^\infty}\oplus\mathbf Q_5/\mathbf Z_5.
\]
The residue $\bar J=t$ is transcendental over $\mathbf F_5$, so
$J$ is transcendental over $\mathbf Q$.

We describe the auxiliary subfield without a hidden transcendence
assumption. The field $k$ is countable, whereas every residue class
in $V$ is uncountable. Enumerate $k$ and choose a representative
$c_n\in V$ of its $n$th element which is transcendental over
$\mathbf Q(J,c_1,\ldots,c_{n-1})$. This is possible because the
algebraic closure of that countable field in $\operatorname{Frac}V$
is countable. Let $C=\mathbf Q(c_1,c_2,\ldots)$. Then $C\cap V\to k$
is onto, $C\cap V$ is dense in $V$, and $J$ is transcendental over $C$.

Since $J$ and $J-1728$ are units, the elliptic scheme
\[
 E_J:\quad Y^2=X^3-\frac{3J}{J-1728}X
                         -\frac{2J}{J-1728}             \tag{C-RD11.3}
\]
has $j$-invariant $J$ and discriminant
$1728^2J^2/(J-1728)^3\in V^\times$.
It is isomorphic over $V$ to $E^{\mathrm{can}}$: elliptic schemes
with the same $j$ away from $0,1728$ differ by a quadratic twist, and
every unit of the complete DVR $V$ is a square because $k$ is
algebraically closed and $2$ is invertible. Fix a compatible BT
isomorphism from $E_J[5^\infty]$ to the split group over $V$.

Each of its generic finite-level isomorphisms is defined over a
finite algebraic extension of $C(J)$ inside $\operatorname{Frac}V$:
the isomorphism scheme of the two finite etale group schemes over
$C(J)$ is finite etale. Let $F$ be the compositum of those countably
many fields. It is algebraic over $C(J)$ and the entire compatible
generic BT isomorphism is defined over $F$. Apply
Lemma~\ref{lem:C-RD11-ring} to this $F$ and the derivation $D(J)=1$,
$D(C)=0$, and obtain $R\subset F$.

Set $S=\operatorname{Spec}R$, $U=D(5)=\operatorname{Spec}F$, and
\[
 M=\mu_{5^\infty,R}\oplus(\mathbf Q_5/\mathbf Z_5)_R,
 \qquad A_U=(\text{equation (C-RD11.3)})_F.
\]
The coefficients of that equation belong to $F$; the chosen finite-level maps give the
required isomorphism on $U$. In fact $J\in R$ and $J-1728$ is a unit
of $R$, so the same equation is already an elliptic scheme over $R$.
Hypothesis a) holds with the identity cover, since $R$ is geometrically
unibranch. Hypothesis b) holds because the punctured strict local
spectrum is connected: a reduced geometrically unibranch local
domain has a reduced irreducible strict henselization, and in this
one-dimensional situation its punctured spectrum has a single
generic point. The closed point has characteristic $5$.

There is already an obstruction modulo $5$; completion is unnecessary
for the contradiction. The homomorphism (C-RD11.1), reduced modulo
$5$, induces
\[
 R/5R\simeq k[\epsilon]/(\epsilon^2),\qquad
 a\longmapsto\bar a+\epsilon\overline{D(a)}.            \tag{C-RD11.4}
\]
The kernel is $5R$ because $a,D(a)\in5V_0$ is equivalent to $a/5\in R$.
For surjectivity use representatives from $C\cap V$ and $z=J-c$,
which maps to $\epsilon$.

Suppose an abelian scheme $A/S$ realized this triple. It has relative
dimension one. Its generic $j$-invariant is $J$, and $R\hookrightarrow F$
is injective, so $j(A)=J\in R$. Under (C-RD11.4) this becomes
$t+\epsilon$. On the other hand, the reduction of $M$ to
$k[\epsilon]/(\epsilon^2)$ is the constant split BT group. If $A_0/k$
is the special fibre of $A$, Serre--Tate applied to the square-zero
thickening $k\hookrightarrow k[\epsilon]/(\epsilon^2)$ identifies
$A$ modulo $5$ with the constant deformation of $A_0$, because their
BT deformations are identified with the same constant group
\cite[Theorem~1.2.1]{Katz1981SerreTate}. Its $j$-invariant must
therefore be $t$. This contradicts $j(A\bmod5)=t+\epsilon$.
\end{proof}

\begin{remark}[Finite symbolic verification]
The exact symbolic certificate \path{verify_nonfinite_reduction.py}
checks the displayed discriminant, $j$-invariant, ordinary Hasse
coefficient and the nonzero dual-number coefficient. It makes no
claim to computationally instantiate the infinite field construction.
\end{remark}

\begin{remark}[Where the infinite system fails to descend]
Over the normalization $V_0$, the chosen generic BT isomorphism
extends, since its finite-level coefficients were chosen inside
$F\cap V$. It does not descend to $R$. The normalization is integral
and radicial but not finite or proper of finite type, so this does
not contradict Proposition~\ref{prop:C-RD1}. It also supplies a
concrete reason that a finite-presentation limit theorem cannot be
applied to an entire BT isomorphism at once.
\end{remark}
